\documentclass[12pt]{article}
\usepackage[utf8]{inputenc}
\usepackage[T1]{fontenc}
\usepackage{hyperref}
\title{Eastern Karaboro Language Resources for K-12 Curriculum Development}
\author{Compiled by Claude Sonnet 5.5 (claude-sonnet-5-5) for Andr\'as Kornai}
\date{}
\begin{document}
\maketitle

\section{Introduction}
Eastern Karaboro (ISO 639-3 xrb; also called Kar, Ker or Kler) is a Senufo language of southern Burkina Faso. Sources cited here give about 35,000 to 40,000 speakers. Its comprehension of Western Karaboro is only partial. Speakers are often literate in French and understand Jula and Bamanankan. This report rests on a single search round, and little material turned up. Nothing found shows that the language is taught in schools.

\section{Primary Educational Resources}
No primers, textbooks or school curricula were found in the search results. A 1994 dictionary was published by a literacy association for Kar, and it may have been used in literacy work. See the next section.

\section{Linguistic and Reference Materials}
\begin{itemize}
\item Wikipedia, Kar language (general overview): \url{https://en.wikipedia.org/wiki/Kar_language}
\item Glottolog reference record for Wichser and B\"uhler 1994, \textit{Dictionnaire Karaboro-Fran\c{c}ais et Fran\c{c}ais-Karaboro, livret 1}, Ouagadougou: Association pour l'Alphab\'etisation en langue Kar and SIL. Only the bibliographic record was seen, not the dictionary itself: \url{https://glottolog.org/resource/reference/id/59899}
\item Wikitongues language page: \url{https://wikitongues.org/languages/xrb/}
\end{itemize}

\section{Religious and Community Resources}
\begin{itemize}
\item Eastern Karaboro New Testament, 1993 edition, Digital Bible Society: \url{https://dbs.org/bibles/XRBWBT}
\item Scripture Earth resources: \url{https://www.scriptureearth.org/00eng.php?iso=xrb}
\item Joshua Project people group profile: \url{https://joshuaproject.net/index.php/people_groups/11696/UV}
\end{itemize}

\section{Digital and Audio Resources}
\begin{itemize}
\item Global Recordings Network audio recordings for Bible teaching: \url{https://www.globalrecordings.net/en/language/xrb}
\end{itemize}

\section{Curriculum Development Context}
Material is thin. It consists of a Latin-script orthography with a New Testament, a literacy-association dictionary, and audio recordings. No school programme was found. A curriculum would probably need to start from the Kar literacy association's earlier work, if that can be located, and to involve the speaker community. The multilingual setting (French, Jula, Bamanankan) should shape the design.

\end{document}